\ifdefined\gkver\else\def\gkver{student}\fi
\documentclass[\gkver]{gaokaozhenti}
\begin{document}

\chapter{2009年四川理综}
%% paperType: 理综物理部分

\begin{enumerate}
\item
%% number: 14
%% paperName: 2009年四川理综第 14 题 6 分
%% typeId: 1
%% type: 单选题
%% chapter: 光学
%% point: 光的电磁说
%% method: 无
%% score: 6
%% degree: 650
%% duplicateId: 0
%% body:
关于电磁波，下列说法正确的是\xzanswer{D}

\fourchoices[answer=D]
{雷达是用 X 光来测定物体位置的设备}
{使电磁波随各种信号而改变的技术叫做解调}
{用红外线照射时，大额钞票上用荧光物质印刷的文字会发出可见光}
{变化的电场可以产生变化的磁场}

%% answer: D
%% memo:
\memoanswer{
雷达是根据超声波测定物体位置的，A 错；使电磁波随各种信号而改变的技术叫做调制，B 错；用紫外线照射时大额钞票上用荧光物质印刷的文字会发出可见光，利用紫外线的荧光效应，C 错；根据麦克斯韦电磁场理论可知变化的电场可以产生变化的磁场、变化的磁场产生电场，D 对。
}

\item
%% number: 15
%% paperName: 2009年四川理综第 15 题 6 分
%% typeId: 1
%% type: 单选题
%% chapter: 曲线运动
%% point: 天体运动
%% method: 无
%% score: 6
%% degree: 450
%% duplicateId: 0
%% body:
据报道，2009 年 4 月 29 日，美国亚利桑那州一天文观测机构发现一颗与太阳系其它行星逆向运行的小行星，代号为 2009HC82。该小行星绕太阳一周的时间为 3.39 年，直径 $2\sim3\Ukm$，其轨道平面与地球轨道平面呈 $155\Udu$ 的倾斜。假定该小行星与地球均以太阳为中心做匀速圆周运动，则小行星和地球绕太阳运动的速度大小的比值为\xzanswer{A}

\fourchoices[answer=A]
{$3.39^{-1/3}$}
{$3.39^{-1/2}$}
{$3.39^{3/2}$}
{$3.39^{2/3}$}

%% answer: A
%% memo:
\memoanswer{
小行星和地球绕太阳作圆周运动，都是由万有引力提供向心力，有 $\frac{GMm}{R^2}=m\left(\frac{2\pi}{T}\right)^2 R$，可知小行星和地球绕太阳运行轨道半径之比为 $R_1:R_2=\sqrt[3]{\frac{T_1^2}{T_2^2}}$，又根据 $v=\sqrt{\frac{GM}{R}}$，联立解得 $v_1:v_2=\sqrt[3]{\frac{T_1}{T_2}}$，已知 $\frac{T_1}{T_2}=\frac{1}{3.39}$，则 $v_1:v_2=\sqrt[3]{\frac{1}{3.39}}$。
}

\item
%% number: 16
%% paperName: 2009年四川理综第 16 题 6 分
%% typeId: 1
%% type: 单选题
%% chapter: 气体性质
%% point: 热力学第一定律
%% method: 无
%% score: 6
%% degree: 650
%% duplicateId: 0
%% body:
关于热力学定律，下列说法正确的是\xzanswer{B}

\fourchoices[answer=B]
{在一定条件下物体的温度可以降到 $0\UK$}
{物体从单一热源吸收的热量可全部用于做功}
{吸收了热量的物体，其内能一定增加}
{压缩气体总能使气体的温度升高}

%% answer: B
%% memo:
\memoanswer{
根据热力学第三定律的绝对零度不可能达到可知 A 错；物体从外界吸收热量、对外做功，根据热力学第一定律可知内能可能增加、减小和不变，C 错；压缩气体，外界对气体作正功，可能向外放热，内能可能减少、温度降低，D 错；物体从单一热源吸收的热量可全部用于做功而引起其他变化是可能的，B 对。
}

\item
%% number: 17
%% paperName: 2009年四川理综第 17 题 6 分
%% typeId: 1
%% type: 单选题
%% chapter: 交流电
%% point: 交流电
%% method: 无
%% score: 6
%% degree: 450
%% duplicateId: 0
%% body:
如图\ref{2009四川理综17a}所示，理想变压器原、副线圈的匝数比为 $10:1$，$R_{1}=20\UO$，$R_{2}=30\UO$，$C$ 为电容器。已知通过 $R_{1}$ 的正弦交流电如图\ref{2009四川理综17b}所示，则\xzanswer{C}

\twopicture[labelA=2009四川理综17a, labelB=2009四川理综17b, numA=甲, numB=乙, w=6.0cm]
{figs/17a.png}
{figs/17b.png}

\fourchoices[answer=C]
{交流电的频率为 $0.02\UHz$}
{原线圈输入电压的最大值为 $200\sqrt{2}\UV$}
{电阻 $R_{2}$ 的电功率约为 $6.67\UW$}
{通过 $R_{3}$ 的电流始终为零}

%% answer: C
%% memo:
\memoanswer{
根据变压器原理可知原副线圈中电流的周期、频率相同，周期为 $0.02\Us$、频率为 $50\UHz$，A 错。由图乙可知通过 $R_{1}$ 的电流最大值为 $I_m=1\UA$、根据欧姆定律可知其最大电压为 $U_m=20\UV$，再根据原副线圈的电压之比等于匝数之比可知原线圈输入电压的最大值为 $200\UV$，B 错；因为电容器有通交流、阻直流的作用，则有电流通过 $R_{3}$ 和电容器，D 错；根据正弦交流电的峰值和有效值关系并联电路特点可知电阻 $R_{2}$ 的电流有效值为 $I=\frac{I_m R_1}{\sqrt{2}R_2}$、电压有效值为 $U=U_m/\sqrt{2}$，电阻 $R_{2}$ 的电功率为 $P_2=UI=\frac{20}{3}\UW$，C 对。
}

\item
%% number: 18
%% paperName: 2009年四川理综第 18 题 6 分
%% typeId: 1
%% type: 单选题
%% chapter: 原子和原子核
%% point: 玻尔模型
%% method: 无
%% score: 6
%% degree: 450
%% duplicateId: 0
%% body:
氢原子能级的示意图如图所示，大量氢原子从 $n=4$ 的能级向 $n=2$ 的能级跃迁时辐射出可见光 a，从 $n=3$ 的能级向 $n=2$ 的能级跃迁时辐射出可见光 b，则\xzanswer{C}

\onepicture[w=4.0cm]{figs/18.png}

\fourchoices[answer=C]
{氢原子从高能级向低能级跃迁时可能会辐射出 $\gamma$ 射线}
{氢原子从 $n=4$ 的能级向 $n=3$ 的能级跃迁时会辐射出紫外线}
{在水中传播时，a 光较 b 光的速度小}
{氢原子在 $n=2$ 的能级时可吸收任意频率的光而发生电离}

%% answer: C
%% memo:
\memoanswer{
$\gamma$ 射线的产生机理是原子核受激发，是原子核变化才产生的，A 错；根据跃迁规律可知高能级向低能级跃迁时辐射光子的能量等于这两个能级差，从 $n=4$ 的能级向 $n=3$ 的能级跃迁时会辐射出的光子能量小于 a 光子的能量、不可能为紫外线，B 错；根据跃迁规律可知从 $n=4$ 向 $n=2$ 跃迁时辐射光子的能量大于从 $n=3$ 向 $n=2$ 跃迁时辐射光子的能量，则可见光 a 的光子能量大于 b，又根据光子能量 $E=h\nu$ 可得 a 光子的频率大于 b，则 a 的折射率大于 b，又 $v=\frac{c}{n}$ 可得在水中传播时，a 光较 b 光的速度小，C 对；欲使在 $n=2$ 能级的氢原子发生电离，吸收的能量一定不小于 $3.4\UeV$，D 错。
}

\item
%% number: 19
%% paperName: 2009年四川理综第 19 题 6 分
%% typeId: 1
%% type: 单选题
%% chapter: 机械波
%% point: 波的图象
%% method: 无
%% score: 6
%% degree: 450
%% duplicateId: 0
%% body:
图示为一列沿 $x$ 轴负方向传播的简谐横波，实线为 $t=0$ 时刻的波形图，虚线为 $t=0.6\Us$ 时的波形图，波的周期 $T>0.6\Us$，则\xzanswer{D}

\onepicture[w=5.5cm]{figs/19.png}

\fourchoices[answer=D]
{波的周期为 $2.4\Us$}
{在 $t=0.9\Us$ 时，P 点沿 $y$ 轴正方向运动}
{经过 $0.4\Us$，P 点经过的路程为 $4\Um$}
{在 $t=0.5\Us$ 时，Q 点到达波峰位置}

%% answer: D
%% memo:
\memoanswer{
根据题意应用平移法可知 $\frac{3}{4}T=0.6\Us$，解得 $T=0.8\Us$，A 错；由图可知振幅 $A=0.2\Um$、波长 $\lambda=8\Um$。$t=0.9\Us=1\frac{1}{8}T$，此时 P 点沿 $y$ 轴负方向运动，B 错；$0.4\Us=\frac{1}{2}T$，运动路程为 $2A=0.4\Um$，C 错；$t=0.5\Us=\frac{5}{8}T=\frac{1}{2}T+\frac{1}{8}T$，波形图中 Q 正在向下振动，从平衡位置向下振动了 $\frac{1}{8}T$，经 $\frac{1}{8}T$ 到波谷，再过 $\frac{1}{2}T$ 到波峰，D 对。
}

\item
%% number: 20
%% paperName: 2009年四川理综第 20 题 6 分
%% typeId: 2
%% type: 多选题
%% chapter: 电场
%% point: 带电粒子的直线运动
%% method: 无
%% score: 6
%% degree: 250
%% duplicateId: 0
%% body:
如图所示，粗糙程度均匀的绝缘斜面下方 O 点处有一正点电荷，带负电的小物体以初速度 $v_{1}$ 从 M 点沿斜面上滑，到达 N 点时速度为零，然后下滑回到 M 点，此时速度为 $v_{2}$（$v_{2}<v_{1}$）。若小物体电荷量保持不变，$OM=ON$，则\xzanswer{AD}

\onepicture[align=right, w=4.2cm]{figs/20.png}

\fourchoices[answer=AD]
{小物体上升的最大高度为 $\frac{v_1^2+v_2^2}{4g}$}
{从 N 到 M 的过程中，小物体的电势能逐渐减小}
{从 M 到 N 的过程中，电场力对小物体先做负功后做正功}
{从 N 到 M 的过程中，小物体受到的摩擦力和电场力均是先增大后减小}

%% answer: AD
%% memo:
\memoanswer{
设斜面倾角为 $\theta$、上升过程沿斜面运动的最大距离为 $L$。因为 $OM=ON$，则 MN 两点电势相等，小物体从 M 到 N、从 N 到 M 电场力做功均为 0。上滑和下滑经过同一个位置时，垂直斜面方向上电场力的分力相等，则经过相等的一小段位移在上滑和下滑过程中电场力分力对应的摩擦力所作的功均为相等的负功，所以上滑和下滑过程克服电场力产生的摩擦力所作的功相等、并设为 $W_1$。在上滑和下滑过程，对小物体，应用动能定理分别有：$-mg\sin\theta L-\mu mg\cos\theta L-W_1=-\frac{mv_1^2}{2}$ 和 $mg\sin\theta L-\mu mg\cos\theta L-W_1=\frac{mv_2^2}{2}$，上两式相减可得 $\sin\theta L=\frac{v_1^2+v_2^2}{4g}$，A 对；由 $OM=ON$，可知电场力对小物体先作正功后作负功，电势能先减小后增大，BC 错；从 N 到 M 的过程中，小物体受到的电场力垂直斜面的分力先增大后减小，而重力分力不变，则摩擦力先增大后减小，在此过程中小物体到 O 的距离先减小后增大，根据库仑定律可知小物体受到的电场力先增大后减小，D 对。
}

\item
%% number: 21
%% paperName: 2009年四川理综第 21 题 6 分
%% typeId: 2
%% type: 多选题
%% chapter: 光学
%% point: 全反射
%% method: 无
%% score: 6
%% degree: 450
%% duplicateId: 0
%% body:
如图所示，空气中有一横截面为半圆环的均匀透明柱体，其内圆半径为 $r$，外圆半径为 $R$，$R=\sqrt{2}r$。现有一束单色光垂直于水平端面 A 射入透明柱体，只经过两次全反射就垂直于水平端面 B 射出。设透明柱体的折射率为 $n$，光在透明柱体内传播的时间为 $t$，若真空中的光速为 $c$，则\xzanswer{AB}

\onepicture[w=4.5cm]{figs/21.png}

\fourchoices[answer=AB]
{$n$ 可能为 $\sqrt{3}$}
{$n$ 可能为 2}
{$t$ 可能为 $\frac{2\sqrt{2}r}{c}$}
{$t$ 可能为 $\frac{4.8r}{c}$}

%% answer: AB
%% memo:
\memoanswer{
只经过两次全反射可知第一次入射角为 $45\Udu$，反射光路图如右图所示。根据全反射可知临界角 $C\leq45\Udu$，再根据 $n=\frac{1}{\sin C}$ 可知 $n\geq\sqrt{2}$；光在透明柱体中运动路程为 $L=4r$，运动时间为 $t=L/v=\frac{4nr}{c}$，则 $t\geq\frac{4\sqrt{2}r}{c}$，CD 均错。

\onepicture[w=4.5cm]{figs/21_an.png}
}

\item
%% number: 22
%% paperName: 2009年四川理综第 22 题 8 分
%% typeId: 4
%% type: 实验题
%% chapter: 运动和力
%% point: 动量守恒定律
%% method: 无
%% score: 8
%% degree: 450
%% duplicateId: 0
%% body:
\begin{enumerate}
	\item
	在弹性限度内，弹簧弹力的大小与弹簧伸长（或缩短）的长度的比值，叫做弹簧的劲度系数。为了测量一轻弹簧的劲度系数，某同学进行了如下实验设计：如图所示，将两平行金属导轨水平固定在竖直向下的匀强磁场中，金属杆 ab 与导轨接触良好，水平放置的轻弹簧一端固定于 O 点，另一端与金属杆连接并保持绝缘。在金属杆滑动的过程中，弹簧与金属杆、金属杆与导轨均保持垂直，弹簧的形变始终在弹性限度内，通过减小金属杆与导轨之间的摩擦和在弹簧形变较大时读数等方法，使摩擦对实验结果的影响可忽略不计。请你按要求帮助该同学解决实验所涉及的两个问题。

	\begin{steps}
		\item
		帮助该同学完成实验设计。请你用低压直流电源（ ）、滑动变阻器（ ）、电流表（ ）、开关（ ）设计一电路图，画在图中虚线框内，并正确连在导轨的 C、D 两端。
		\item
		若已知导轨间的距离为 $d$，匀强磁场的磁感应强度为 $B$，正确连接电路后，闭合开关，使金属杆随挡板缓慢移动，当移开挡板且金属杆静止时，测出通过金属杆的电流为 $I_{1}$，记下金属杆的位置，断开开关，测出弹簧对应的长度为 $x_{1}$；改变滑动变阻器的阻值，再次让金属杆静止时，测出通过金属杆的电流为 $I_{2}$，弹簧对应的长度为 $x_{2}$，则弹簧的劲度系数 $k=$\tkanswer[2cm]{}。
	\end{steps}
	\item
	气垫导轨（如图\ref{2009四川理综22f}）工作时，空气从导轨表面的小孔喷出，在导轨表面和滑块内表面之间形成一层薄薄的空气层，使滑块不与导轨表面直接接触，大大减小了滑块运动时的阻力。为了验证动量守恒定律，在水平气垫导轨上放置两个质量均为 $a$ 的滑块，每个滑块的一端分别与穿过打点计时器的纸带相连，两个打点计时器所用电源的频率均为 $b$。气垫导轨正常工作后，接通两个打点计时器的电源，并让两滑块以不同的速度相向运动，两滑块相碰后粘在一起继续运动。图\ref{2009四川理综22g}为某次实验打出的、点迹清晰的纸带的一部分，在纸带上以同间距的 6 个连续点为一段划分纸带，用刻度尺分别量出其长度 $s_{1}$、$s_{2}$ 和 $s_{3}$。若题中各物理量的单位均为国际单位，那么，碰撞前两滑块的动量大小分别为 \tkanswer[2cm]{}、\tkanswer[2cm]{}，两滑块的总动量大小为 \tkanswer[2cm]{}；碰撞后两滑块的总动量大小为 \tkanswer[2cm]{}。重复上述实验，多做几次。若碰撞前、后两滑块的总动量在实验误差允许的范围内相等，则动量守恒定律得到验证。
\end{enumerate}

\onepicture[align=right, w=5.0cm]{figs/22a.png}

\fourpicture[capA=低压直流电源, capB=滑动变阻器, capC=电流表, capD=开关, showprefix=false, wA=1.0cm, wB=1.8cm, wC=1.0cm, wD=1.4cm, gap=1.6cm]
{figs/22b.png}
{figs/22c.png}
{figs/22d.png}
{figs/22e.png}

\twopicture[labelA=2009四川理综22f, labelB=2009四川理综22g, numA=甲, numB=乙, w=6.5cm]
{figs/22f.png}
{figs/22g.png}

%% answer:
\jdanswer{
\begin{enumerate}
	\item
	①低压直流电源 $E$、滑动变阻器 $R$、电流表、开关 $S$ 串接在 CD 两点之间，电路图如解析图所示。

	②设弹簧原长为 $L_0$，应用胡克定律有 $k(x_1-L_0)=BI_1d$、$k(x_2-L_0)=BI_2d$，两式相减可得 $k(x_1-x_2)=B(I_1-I_2)d$，解得 $k=\frac{Bd(I_1-I_2)}{x_1-x_2}$。
	\item
	$0.2ab s_{3}$，$0.2ab s_{1}$（两空可互换），$0.2ab(s_1-s_3)$，$0.4ab s_2$
\end{enumerate}
}

%% memo:
\memoanswer{
（1）①低压直流电源 $E$、滑动变阻器 $R$、电流表、开关 $S$ 串接在 CD 两点之间，如图所示。

\onepicture[w=6.0cm]{figs/22_an.png}

②设弹簧原长为 $L_0$，应用胡克定律有 $k(x_1-L_0)=BI_1d$、$k(x_2-L_0)=BI_2d$，两式相减可得 $k(x_1-x_2)=B(I_1-I_2)d$，解得 $k=\frac{Bd(I_1-I_2)}{x_1-x_2}$；法二：根据胡克定律 $F=kx$ 可得 $\Delta F=k\Delta x$，则 $k=\frac{\Delta F}{\Delta x}=\frac{Bd(I_1-I_2)}{x_1-x_2}$。

（2）动量 $P=mv$，根据 $v=s/(5T)$ 可知两滑块碰前的速度分别为 $v_1=0.2s_1b$、$v_2=0.2s_3b$，则碰前动量分别为 $0.2abs_1$ 和 $0.2abs_3$，总动量大小为 $av_1-av_2=0.2ab(s_1-s_3)$；碰撞后两滑块的总动量大小为 $2av=2a s_2/(5T)=0.4abs_2$。
}

\item
%% number: 23
%% paperName: 2009年四川理综第 23 题 16 分
%% typeId: 6
%% type: 计算题
%% chapter: 功和能
%% point: 交通工具启动
%% method: 无
%% score: 16
%% degree: 450
%% duplicateId: 0
%% body:
图示为修建高层建筑常用的塔式起重机。在起重机将质量 $m=5\times10^{3}\Ukg$ 的重物竖直吊起的过程中，重物由静止开始向上作匀加速直线运动，加速度 $a=0.2\Umsq$，当起重机输出功率达到其允许的最大值时，保持该功率直到重物做 $v_{m}=1.02\Ums$ 的匀速运动。取 $g=10\Umsq$，不计额外功。求：

\begin{enumerate}
	\item
	起重机允许输出的最大功率。
	\item
	重物做匀加速运动所经历的时间和起重机在第 2 秒末的输出功率。
\end{enumerate}

\onepicture[align=right, w=4.5cm]{figs/23.png}

%% answer:
\jdanswer{
\begin{enumerate}
	\item
	$P_{0}=5.1\times10^{4}\UW$
	\item
	$t_{1}=5\Us$，$P=2.04\times10^{4}\UW$
\end{enumerate}
}

%% memo:
\memoanswer{
（1）设起重机允许输出的最大功率为 $P_0$，重物达到最大速度时，拉力 $F_0$ 等于重力。

$P_0=F_0v_m$ ①

$P_0=mg$ ②

代入数据，有：$P_0=5.1\times10^{4}\UW$ ③

（2）匀加速运动结束时，起重机达到允许输出的最大功率，设此时重物受到的拉力为 $F$，速度为 $v_1$，匀加速运动经历时间为 $t_1$，有：

$P_0=F_0v_1$ ④

$F-mg=ma$ ⑤

$v_1=at_1$ ⑥

由③④⑤⑥，代入数据，得：$t_1=5\Us$ ⑦

$t=2\Us$ 时，重物处于匀加速运动阶段，设此时速度为 $v_2$，输出功率为 $P$，则

$v_2=at$ ⑧

$P=Fv_2$ ⑨

由⑤⑧⑨，代入数据，得：$P=2.04\times10^{4}\UW$。
}

\item
%% number: 24
%% paperName: 2009年四川理综第 24 题 19 分
%% typeId: 6
%% type: 计算题
%% chapter: 电场
%% point: 带电粒子的圆周运动
%% method: 无
%% score: 19
%% degree: 450
%% duplicateId: 0
%% body:
如图所示，直线形挡板 $p_1p_2p_3$ 与半径为 $r$ 的圆弧形挡板 $p_3p_4p_5$ 平滑连接并安装在水平台面 $b_1b_2b_3b_4$ 上，挡板与台面均固定不动。线圈 $c_1c_2c_3$ 的匝数为 $n$，其端点 $c_1$、$c_3$ 通过导线分别与电阻 $R_1$ 和平行板电容器相连，电容器两极板间的距离为 $d$，电阻 $R_1$ 的阻值是线圈 $c_1c_2c_3$ 阻值的 2 倍，其余电阻不计，线圈 $c_1c_2c_3$ 内有一面积为 $S$、方向垂直于线圈平面向上的匀强磁场，磁场的磁感应强度 $B$ 随时间均匀增大。质量为 $m$ 的小滑块带正电，电荷量始终保持为 $q$，在水平台面上以初速度 $v_0$ 从 $p_1$ 位置出发，沿挡板运动并通过 $p_5$ 位置。若电容器两板间的电场为匀强电场，$p_1$、$p_2$ 在电场外，间距为 $L$，其间小滑块与台面的动摩擦因数为 $\mu$，其余部分的摩擦不计，重力加速度为 $g$。求：

\begin{enumerate}
	\item
	小滑块通过 $p_{2}$ 位置时的速度大小。
	\item
	电容器两极板间电场强度的取值范围。
	\item
	经过时间 $t$，磁感应强度变化量的取值范围。
\end{enumerate}

\onepicture[align=right, w=7.5cm]{figs/24.png}

%% answer:
\jdanswer{
\begin{enumerate}
	\item
	$\sqrt{v_0^2-2\mu gL}$
	\item
	$0<E\leq\frac{m(v_0^2-2\mu gL)}{5qr}$
	\item
	$0<\Delta B\leq\frac{3md(v_0^2-2\mu gL)}{10nSqr}t$
\end{enumerate}
}

%% memo:
\memoanswer{
（1）小滑块运动到位置 $p_2$ 时速度为 $v_1$，由动能定理有：

$-\mu mgL=\frac{1}{2}mv_1^2-\frac{1}{2}mv_0^2$ ①

$v_1=\sqrt{v_0^2-2\mu gL}$ ②

（2）由题意可知，电场方向如图，若小滑块能通过位置 $p$，则小滑块可沿挡板运动且通过位置 $p_5$，设小滑块在位置 $p$ 的速度为 $v$，受到的挡板的弹力为 $N$，匀强电场的电场强度为 $E$，由动能定理有：

$-\mu mgL-2rEqs=\frac{1}{2}mv_1^2-\frac{1}{2}mv_0^2$ ③

当滑块在位置 $p$ 时，由牛顿第二定律有：$N+Eq=m\frac{v^2}{r}$ ④

由题意有：$N\geq0$ ⑤

由以上三式可得：$E\leq\frac{m(v_0^2-2\mu gL)}{5qr}$ ⑥

$E$ 的取值范围：$0<E\leq\frac{m(v_0^2-2\mu gL)}{5qr}$ ⑦

\onepicture[w=4.0cm]{figs/24_an.png}

（3）设线圈产生的电动势为 $E_1$，其电阻为 $R$，平行板电容器两端的电压为 $U$，$t$ 时间内磁感应强度的变化量为 $\Delta B$，得：

$U=Ed$ ⑧

由法拉第电磁感应定律得 $E_1=n\frac{\Delta BS}{t}$ ⑨

由全电路的欧姆定律得 $E_1=I(R+2R)$ ⑩

$U=2RI$

经过时间 $t$，磁感应强度变化量的取值范围：$0<\Delta B\leq\frac{3md(v_0^2-2\mu gL)}{10nSqr}t$。
}

\item
%% number: 25
%% paperName: 2009年四川理综第 25 题 20 分
%% typeId: 6
%% type: 计算题
%% chapter: 磁场
%% point: 带电粒子在复合场中的运动
%% method: 无
%% score: 20
%% degree: 250
%% duplicateId: 0
%% body:
如图所示，轻弹簧一端连于固定点 O，可在竖直平面内自由转动，另一端连接一带电小球 P，其质量 $m=2\times10^{-2}\Ukg$，电荷量 $q=0.2\UC$。将弹簧拉至水平后，以初速度 $v_{0}=20\Ums$ 竖直向下射出小球 P，小球 P 到达 O 点的正下方 $O_1$ 点时速度恰好水平，其大小 $v=15\Ums$。若 O、$O_1$ 相距 $R=1.5\Um$，小球 P 在 $O_1$ 点与另一由细绳悬挂的、不带电的、质量 $M=1.6\times10^{-1}\Ukg$ 的静止绝缘小球 N 相碰。碰后瞬间，小球 P 脱离弹簧，小球 N 脱离细绳，同时在空间加上竖直向上的匀强电场 E 和垂直于纸面的磁感应强度 $B=1\UT$ 的弱强磁场。此后，小球 P 在竖直平面内做半径 $r=0.5\Um$ 的圆周运动。小球 P、N 均可视为质点，小球 P 的电荷量保持不变，不计空气阻力，取 $g=10\Umsq$。那么，

\begin{enumerate}
	\item
	弹簧从水平摆至竖直位置的过程中，其弹力做功为多少？
	\item
	请通过计算并比较相关物理量，判断小球 P、N 碰撞后能否在某一时刻具有相同的速度。
	\item
	若题中各量为变量，在保证小球 P、N 碰撞后某一时刻具有相同速度的前提下，请推导出 $r$ 的表达式（要求用 $B$、$q$、$m$、$\theta$ 表示，其中 $\theta$ 为小球 N 的运动速度与水平方向的夹角）。
\end{enumerate}

\onepicture[align=right, w=4.5cm]{figs/25.png}

%% answer:
\jdanswer{
\begin{enumerate}
	\item
	$-2.05\UJ$
	\item
	P、N 的速度在同一时刻不可能相同
	\item
	$r=\frac{[(2n+1)\pi+\theta]m^{2}g}{B^{2}q^{2}\sin\theta}$（$n=0,1,2\ldots$）
\end{enumerate}
}

%% memo:
\memoanswer{
（1）设弹簧的弹力做功为 $W$，有：

$mgR+W=\frac{1}{2}mv^2-\frac{1}{2}mv_0^2$ ①

代入数据，得：$W=-2.05\UJ$ ②

（2）由题给条件知，N 碰后作平抛运动，P 所受电场力和重力平衡，P 带正电荷。设 P、N 碰后的速度大小分别为 $v_1$ 和 $V$，并令水平向右为正方向，有：

$mv=\pm mv_1+MV$ ③

而：

$v_1=\frac{Bqr}{m}$ ④

若 P、N 碰后速度同向时，计算可得 $V<v_1$，这种碰撞不能实现。P、N 碰后瞬时必为反向运动。有：$V=\frac{mv+Bqr}{M}$ ⑤

P、N 速度相同时，N 经过的时间为 $t_N$，P 经过的时间为 $t_P$。设此时 N 的速度 $V_1$ 的方向与水平方向的夹角为 $\theta$，有：

$\cos\theta=\frac{V}{V_1}=\frac{V}{v_1}$ ⑥

$gt_N=V_1\sin\theta=v_1\sin\theta$ ⑦

代入数据，得：$t_N=\frac{\sqrt{3}}{4}\Us$ ⑧

对小球 P，其圆周运动的周期为 $T$，有：

$T=\frac{2\pi m}{Bq}$ ⑨

经计算得：$t_N<T$，

P 经过 $t_P$ 时，对应的圆心角为 $\alpha$，有：$t_P=\frac{\alpha}{2\pi}T$ ⑩

\onepicture[w=5.5cm]{figs/25_an.png}

当 B 的方向垂直纸面朝外时，P、N 的速度相同，如图可知，有：$\alpha_1=\pi+\theta$

联立相关方程得：$t_{P1}=\frac{2\pi}{15}\Us$

比较得，$t_N\neq t_{P1}$，在此情况下，P、N 的速度在同一时刻不可能相同。

当 B 的方向垂直纸面朝里时，P、N 的速度相同，同样由图，有：$\alpha_2=\pi-\theta$，

同上得：$t_{P2}=\frac{\pi}{15}\Us$，

比较得，$t_N\neq t_{P2}$，在此情况下，P、N 的速度在同一时刻也不可能相同。

（3）当 B 的方向垂直纸面朝外时，设在 $t$ 时刻 P、N 的速度相同，$t_N=t_P=t$，

再联立④⑦⑨⑩解得：$r=\frac{[(2n+1)\pi+\theta]m^{2}g}{B^{2}q^{2}\sin\theta}$（$n=0,1,2\ldots$）

当 B 的方向垂直纸面朝里时，设在 $t$ 时刻 P、N 的速度相同 $t_N=t_P=t$，

同理得：$r=\frac{(\pi-\theta)m^{2}g}{B^{2}q^{2}\sin\theta}$，

考虑圆周运动的周期性，有：$r=\frac{[(2n+1)\pi+\theta]m^{2}g}{B^{2}q^{2}\sin\theta}$（$n=0,1,2\cdots$）（给定的 $B$、$q$、$r$、$m$、$\theta$ 等物理量决定 $n$ 的取值）。
}

\end{enumerate}

\end{document}
